Die sorgfältige Arbeit mit Literatur ist in allen Wissenschaftszweigen die
Basis für jede Veröffentlichung. Hierbei ist es nicht nur mit der Erstellung
eines Literaturverzeichnisses getan. Die eigentliche Arbeit beginnt bereits
bevor man mit dem Schreiben der eigenen Veröffentlichung beginnt. Folgende
Arbeitsschritte sind für die Erstellung der Bibliografien notwendig.
%
\begin{compactenum}
\item Literaturrecherche und Erstellung einer bzw. mehrerer
      Literaturdatenbanken; Dieser Schritt ist unabhängig vom \LaTeX-Workflow.
\item Einbindung der Litaturdatenbanken in das \LaTeX-Dokument
\item Erstellung von Literaturverweisen im \LaTeX-Dokument
\item Ggf. Anpassung der Gestaltung des Literaturverzeichnisses
\end{compactenum}

\LaTeX\ bietet vielfältige Wege eine Bibliografie zu erstellen. Im einfachsten
Fall wird die Bibliografie manuell durch die \texttt{thebibliography}-Umgebung
erstellt. Dieser Ansatz führt schnell zu Ergebnissen und ist für kleinere
Seminararbeiten durchaus praktikabel. Für umfangreichere Arbeit nimmt der
Aufwand hierfür aber schnell überhand. Im Rahmen dieser Vorlag benutzen wir
deshalb das Zusatzpaket~\texttt{biblatex}, welches folgende Vorteile bietet.
%
\begin{compactitem}
\item Eine einmal erstellte Literaturdatenbank kann in späteren
      Veröffentlichungen wiederverwendet werden.
\item Literaturmarken werden automatisch anhand vorgegebener Schemata erstellt.
\item Die Formatierung des Literaturverzeichnisses erfolgt einheitlich gemäß
      einschlägiger Vorgaben.
\item Es werden stets nur die Literaturstellen aus der Datenbank in das
      Verzeichnis übernommen, auf die auch im Dokument verwiesen wird.
\end{compactitem}

Die Gestaltung von Bibliografien ist ein komplexes Thema. Im Rahmen dieser
Anleitung werden wir nur darauf eingehen, wie sich das Literaturverzeichnis
technisch in die Vorlage einfügt. Für Detailfragen oder weitergehende
Informationen zur Gestaltung sei auf \cite{voss_bibliografien} und
\cite[Kapitel~13]{voss_wissenschaftliche_arbeit} verwiesen. Das
Paket~\texttt{biblatex} verfügt über vielfältige Einstellmöglichkeiten, die im
Handbuch \cite{pkgmanual_biblatex} beschrieben sind. Eine Kurzübersicht zum
Paket und zum Aufbau der Literaturdatenbank findet sich unter
\cite{rees_biblatex_cheat}.
